%% Chapter 6 — Section 6.3: Langevin Equation and its Numerical Solution
\section{Langevin Equation and its Numerical Solution}
\label{sec:6.3}

In Section~\ref{sec:6.1} we introduced ULA by analogy to the gradient descent optimization
algorithm. Another insightful interpretation of ULA is to view it as a discretization of the
Langevin stochastic differential equation. In fact, this perspective predates the optimization
one and explains the name of the algorithm. Here, we summarize some key ideas without delving
into the important technical issues that are needed for a rigorous treatment of this subject.

Given a potential $V : \R^d \to \R$, the (overdamped) Langevin equation is given by
\begin{equation}
  \label{eq:6.7}
  dX(t) = -\nabla V\!\bigl(X(t)\bigr)\,dt + \sqrt{2}\,dW(t), \quad 0 \le t < \infty,
\end{equation}
where $W$ denotes a standard Brownian motion in $\R^d$. The solution to \eqref{eq:6.7} is a
continuous-time stochastic process $\{X(t)\}_{t \ge 0}$. Given a final time $T > 0$, a
\emph{trajectory} of this process is a random function
\[
  [0, T] \to \R^d, \quad t \mapsto X(t).
\]
Figure~\ref{fig:6.2} shows a trajectory of Langevin dynamics.

\begin{figure}[htbp]
  \centering
  \includegraphics[width=0.75\textwidth]{figures/fig_06_2}
  \caption{\textit{Random trajectory of Langevin equation with the potential
    $V(x) = x^4 - x^2 - 0.4x$ shown in Figure~\ref{fig:6.1}. The process spends most of its
    time around the modes of the invariant distribution $f \propto \exp(-V)$.}}
  \label{fig:6.2}
\end{figure}

The p.d.f.\ $\pi_t(x) = \pi(x, t)$ of $X(t)$ defined by \eqref{eq:6.7} satisfies the
Fokker--Planck equation
\begin{align*}
  \frac{\partial \pi}{\partial t} &= \operatorname{div}(\nabla V\,\pi + \nabla\pi), \\
  \pi(x, 0) &= \pi_0(x).
\end{align*}
Notice that for $f \propto \exp(-V)$, we have that $\nabla V = -\nabla\log f = \frac{-\nabla f}{f}$,
which implies that
\[
  \operatorname{div}(\nabla V f + \nabla f) = 0.
\]
Thus, $f$ is an invariant distribution for the Fokker--Planck equation, which means that if
$X(0) \sim f$, then $X(t) \sim f$ for all $t \ge 0$. Moreover, the key property of Langevin
equation that makes it appealing for sampling is that, under mild asymptotic growth conditions
on the potential $V$, it defines an ergodic process. This means that, for large $t$, $X(t)$ is
approximately distributed like $f$, regardless of the initial distribution $\pi_0$ of $X(0)$.
An illustration is given in Figure~\ref{fig:6.3}.

\begin{figure}[htbp]
  \centering
  \includegraphics[width=0.75\textwidth]{figures/fig_06_3}
  \caption{\textit{Histograms of the distribution $\pi_t$ of $X(t)$ for Langevin equation with
    the potential $V(x) = x^4 - x^2 - 0.4x$ in Figure~\ref{fig:6.1}. For large $t$, $\pi_t$
    is close to the target density $f \propto \exp(-V)$.}}
  \label{fig:6.3}
\end{figure}

Although any given trajectory may spend large amounts of time hovering around some local minima
of $V$, if $T$ is large enough and $A \subset E = \R^d$, the proportion of time that the
trajectory spends in $A$ will approximately be $\int_A f(x)\,dx$. Such sample path ergodicity
inspires a natural way to estimate expectations with respect to $f$:
\begin{equation}
  \label{eq:6.8}
  I_f[h] := \int_E h(x)f(x)\,dx \approx \frac{1}{T}\int_0^T h\!\bigl(X(t)\bigr)\,dt,
\end{equation}
where $h$ is a given test function. The quality of the approximation of the spatial average via
the temporal average along a trajectory improves as the end-time $T$ increases.

In practice, Langevin equation \eqref{eq:6.7} can only be solved analytically in simple cases,
e.g.\ for quadratic potentials $V$ that result in linear stochastic differential equations.
Outside these cases, the Euler--Maruyama method is a standard approach to find numerical
solutions. Let $0 < t_1 < \cdots < t_N = T$ with $t_n - t_{n-1} = \epsilon$, $n = 1,\ldots,N$.
Applied to the Langevin equation \eqref{eq:6.7}, Euler--Maruyama approximates $X(t_n)$ with
$X^{(n)}$, where
\[
  X^{(n+1)} = X^{(n)} - \epsilon\nabla V\!\bigl(X^{(n)}\bigr) + \sqrt{2\epsilon}\,\xi^{(n)},
  \qquad \xi^{(n)} \iid \Normal(0, I).
\]
In other words, ULA agrees with an Euler--Maruyama discretization of Langevin equation.
In analogy with \eqref{eq:6.8}, one can then approximate expectations with respect to the
target using the discretized process
\[
  I_f[h] := \int_E h(x)f(x)\,dx \approx \frac{1}{N}\sum_{n=1}^N h\!\bigl(X^{(n)}\bigr).
\]
An important observation is that Euler--Maruyama does not preserve the ergodic property of
Langevin equation, i.e.\ the invariant distribution for the chain $\{X^{(n)}\}_{n \ge 1}$
will no longer be $f$. However, we can use the Markov kernel
\[
  q_{\mathrm{ULA}}(x,\,\cdot\,) = \Normal\!\bigl(x + \epsilon\nabla\log f(x),\, 2\epsilon I\bigr)
\]
implicitly defined by ULA as a proposal kernel within a Metropolis Hastings algorithm. The use
of this proposal kernel along with its correction via the standard Metropolis Hastings acceptance
probability
\[
  a(x, z) := \min\!\left\{1,\; \frac{f(z)}{f(x)}\,
    \frac{q_{\mathrm{ULA}}(z, x)}{q_{\mathrm{ULA}}(x, z)}\right\}
\]
defines MALA. In this light, the Metropolis Hastings accept/reject step in MALA can be
interpreted as removing the bias introduced by Euler--Maruyama discretization of Langevin
equation.
